
\section{Le pont vers le papier principal : décomposition par blocs}
\label{sec:blocks}

Cette section relie la théorie qui précède à l'indicateur du papier principal.
Elle établit l'identité sur laquelle repose la \og carte de contagion \fg{} et
montre où le canal de dépendance se place dans l'indice structurel.

\subsection{L'identité}

Partitionnons les $N$ séries en $B$ blocs (par exemple les segments du CISS :
monétaire, obligataire, actions, intermédiaires financiers, change), et notons
$R_{bb}$ les blocs diagonaux de la matrice de corrélation $R$ et
$R_D=\operatorname{diag}(R_{11},\dots,R_{BB})$.

\begin{proposition}[Décomposition intra/inter-blocs]
\label{prop:blocks}
\begin{equation}
\label{eq:blocks}
\underbrace{-\tfrac12\log\det R}_{\text{dépendance totale}}
\;=\;\sum_{b=1}^{B}\underbrace{\left(-\tfrac12\log\det R_{bb}\right)}_{\text{intra-bloc }b}
\;+\;\underbrace{\left(-\tfrac12\log\det\big(R_D^{-1/2}RR_D^{-1/2}\big)\right)}_{\text{inter-blocs}} ,
\end{equation}
et chacun des $B+1$ termes est positif ou nul. Le terme inter-blocs est la
multi-information gaussienne entre les blocs,
$\KL{\mathcal N(0,R)}{\bigotimes_b\mathcal N(0,R_{bb})}$ ; il est nul si et
seulement si les blocs sont deux à deux non corrélés.
\end{proposition}

\begin{proof}
$R=R_D^{1/2}\,\big(R_D^{-1/2}RR_D^{-1/2}\big)\,R_D^{1/2}$ et le déterminant est
multiplicatif, donc $\log\det R=\log\det R_D+\log\det(R_D^{-1/2}RR_D^{-1/2})$, avec
$\log\det R_D=\sum_b\log\det R_{bb}$. Chaque $R_{bb}$ est une matrice de
corrélation, donc $-\tfrac12\log\det R_{bb}\ge0$ par l'inégalité de Hadamard
(annexe~\ref{app:gauss}). La matrice $\tilde R=R_D^{-1/2}RR_D^{-1/2}$ a des blocs
diagonaux égaux à l'identité, donc une diagonale unité, et Hadamard s'applique
encore. Enfin la proposition~\ref{prop:klgauss} avec $\mu_1=\mu_0$,
$\Sigma_1=R$, $\Sigma_0=R_D$ donne
$\tfrac12[\operatorname{tr}(R_D^{-1}R)-N-\log\det(R_D^{-1}R)]$, et
$\operatorname{tr}(R_D^{-1}R)=\sum_b\operatorname{tr}(I)=N$.
\end{proof}

\paragraph{Deux blocs : les corrélations canoniques.} Pour $B=2$, en notant
$C=R_{11}^{-1/2}R_{12}R_{22}^{-1/2}$ et $\rho_1,\dots,\rho_k$ ses valeurs
singulières (les corrélations canoniques entre les deux blocs),
$\det\tilde R=\det(I-CC^\top)=\prod_j(1-\rho_j^2)$ et le terme inter-blocs vaut
$-\tfrac12\sum_j\log(1-\rho_j^2)$. C'est l'information mutuelle gaussienne
entre les deux segments, somme des informations portées par chaque paire
canonique. Pour $B>2$, l'identité~\eqref{eq:blocks} s'applique de façon
emboîtée et permet d'attribuer la contagion à des paires ou des groupes de
segments par règle de chaînage.

\subsection{Où se place le canal de dépendance dans l'indice structurel}

La règle de chaînage de l'annexe~\ref{app:index} écrit
$\JJ=\sum_i\Dun(x_i)+I(x)$. En remplaçant la multi-information $I(x)$ par sa
part gaussienne $-H^{dep}$ et en appliquant~\eqref{eq:blocks} :
\begin{equation}
\label{eq:jblocks}
\JJ\;\approx\;\underbrace{\sum_i\Dun(x_i)}_{\text{forme}}
\;+\;\underbrace{\sum_b\left(-\tfrac12\log\det R_{bb}\right)}_{\text{intra}}
\;+\;\underbrace{\left(-\tfrac12\log\det\tilde R\right)}_{\text{inter}} .
\end{equation}
Le papier principal conserve les deux derniers termes et abandonne le premier,
qui est celui que le témoin reproduit (section~\ref{sec:program}). Le CISS
entropique est donc l'indice structurel amputé de son canal de forme, et sa
décomposition est exacte, sans approximation, parce qu'elle ne porte que sur la
partie gaussienne.

\subsection{Trois propriétés utiles à l'indicateur}

\paragraph{Invariance d'échelle.} Les termes de~\eqref{eq:blocks} ne dépendent
que de $R$ : ils sont invariants à toute remise à l'échelle série par série,
ce qui permet d'agréger des marchés aux volatilités très différentes sans
pondération ad hoc. La dévolatilisation de \citet{forbes2002} n'agit donc pas
sur la formule mais sur le biais de l'\emph{estimateur} de $R$ en période
d'hétéroscédasticité.

\paragraph{Version en rangs.} En remplaçant $R$ par la corrélation des scores
normaux $\Phi^{-1}(\hat F_i(x_i))$, on estime la multi-information d'une copule
gaussienne. La quantité devient invariante à toute transformation monotone des
marges, ce qui explique sa robustesse aux queues observée en phase~A.

\paragraph{Normalisation en $N$.} Pour une corrélation constante $\rho$,
$\det R=(1-\rho)^{N-1}\big(1+(N-1)\rho\big)$, de sorte que $-H^{dep}$ croît à peu
près linéairement en $N$. Diviser par $N-1$, ou rapporter la valeur à la
corrélation constante équivalente $\rho^\star$ qui donne le même déterminant,
rend l'indicateur comparable d'un panel à l'autre. Pour $\rho^\star$ petit,
$-H^{dep}\approx\tfrac14N(N-1)\rho^{\star2}$.

\section{Indicateurs de suivi et solvabilité en assurance}
\label{sec:monitoring}

\paragraph{Deux indicateurs de la v1.} La séparation saut/diffusion livrait deux
quantités quotidiennes : la composante de saut positif standardisée de $\JJ$
(quelle quantité d'information est arrivée aujourd'hui) et $\JJ_t$ rapporté à
son maximum glissant (quelle part n'est pas encore résorbée). Ni l'un ni l'autre
n'améliore la prévision de la volatilité (H7), et tous deux reposent sur des
statistiques que le témoin n'identifie pas. Ils ne sont pas repris.

\paragraph{Assurance.} Trois liens restent valables en principe. Un scénario en
boule d'entropie produit une loi stressée complète, objet naturel pour l'ORSA.
La décomposition des bundles classiques concerne la formule standard de
Solvabilité~II, dont le crédit de diversification repose sur une matrice de
corrélation fixe. Enfin les mesures de risque entropiques sont cohérentes
\citep{ahmadijavid2012,artzner1999}. Rien n'a cependant été testé sur des
passifs, et les panels courts et peu fréquents des données d'assurance mettraient
à mal chacun des estimateurs employés.

\section{Questions ouvertes}
\label{sec:open}

\begin{enumerate}
\item \textbf{Arrivées auto-excitantes.} Le regroupement des sauts observé dans
l'étude d'événement de la v1 contredit la mesure de Poisson de~\eqref{eq:ou}.
Un processus de Hawkes est la spécification suivante.
\item \textbf{Positivité.} $\JJ\ge0$ alors qu'un Ornstein--Uhlenbeck à terme
brownien peut devenir négatif. Il faudrait soit supprimer $\sigma_\JJ$ (forme
pure de \citealt{bns2001}), soit modéliser $\log\JJ$.
\item \textbf{Dépendance de queue non linéaire.} $\JJ$ et le canal de dépendance
ne voient que la part gaussienne de la multi-information. Une estimation de la
multi-information de copule complète (par exemple par plus proches voisins sur
les rangs) mesurerait ce qui manque, au prix d'une variance d'estimation bien
plus forte.
\item \textbf{Queue du flux d'information.} Au-delà d'une vingtaine d'années,
les quantiles de $I^{(h)}$ sont extrapolés. Un ajustement de Pareto généralisée
serait la correction disciplinée.
\item \textbf{Relèvement minimal et scénarios publiés.} La
proposition~\ref{prop:lift} donne un étalon : pour tout scénario publié, le
rapport entre son prix ambiant et le prix minimal $g(u,v)$ de sa projection
mesure combien d'information il dépense hors de la direction qui compte pour le
portefeuille. C'est un prolongement direct de l'audit du papier principal.
\end{enumerate}
